\section*{Research Foundations and Further Reading}
\addcontentsline{toc}{section}{Research Foundations and Further Reading}

The mechanisms analyzed throughout this chapter build upon foundational scholarship in industrial organization, competitive strategy, and customer decision theory:

\begin{itemize}[leftmargin=14pt,itemsep=3pt,topsep=2pt]
  \item \textbf{[R1] Michael E. Porter (2008), ``The Five Competitive Forces That Shape Strategy,''} \textit{Harvard Business Review.} Foundational analysis of industry rivalry, supplier leverage, buyer power, substitution, and barriers to entry as structural determinants of economic surplus.
  \item \textbf{[R2] Adam Brandenburger and Barry Nalebuff (1996), \textit{Co-opetition}}, and subsequent work on complementor dynamics. Demonstrates how complementary products shape aggregate demand and bilateral bargaining.
  \item \textbf{[R3] Orit Gadiesh and James L. Gilbert (1998), ``How to Map Your Industry's Profit Pool,''} \textit{Harvard Business Review.} Establishes the operational framework for identifying where operating profits actually accumulate across an industry value chain.
  \item \textbf{[R4] Frank M. Bass (1969), ``A New Product Growth Model for Consumer Durables,''} \textit{Management Science.} The classical mathematical model for technology diffusion and market adoption curves.
  \item \textbf{[R5] Marvin B. Lieberman and David B. Montgomery (1988), ``First-Mover Advantages,''} \textit{Strategic Management Journal.} Rigorous empirical review of the conditions under which early market entry confers durable advantages versus fatal liabilities.
  \item \textbf{[R6] Clayton Christensen (1997), \textit{The Innovator's Dilemma}}, and the Christensen Institute's Disruptive Innovation Theory. Details the mechanisms through which simple, lower-cost offerings capture overserved customer tiers.
  \item \textbf{[R7] Ministry of Education, Government of India (8 July 2026), ``AISHE 2022-23 and 2023-24 Reports Released.''} Official government census on national higher-education enrollment, institution counts, and STEM demographics.
  \item \textbf{[R8] NASSCOM (19 August 2024), ``Advancing India's AI Skills: Interventions and Programmes Needed.''} Industry research on software enterprise AI deployment and technical labor shortages.
  \item \textbf{[R9] Stacey Frederick (2019), ``Global Value Chain Mapping,''} \textit{Handbook on Global Value Chains.} Practical methodologies for decomposing industry activities, trade flows, and governance structures.
  \item \textbf{[R10] Michael E. Porter (1996), ``What Is Strategy?''} \textit{Harvard Business Review.} Explains why sustainable competitive advantage requires explicit operational trade-offs and mutually reinforcing activities.
  \item \textbf{[R11] Clayton Christensen et al. (2016), ``Know Your Customers' Jobs to Be Done,''} \textit{Harvard Business Review.} Details the causal framework of customer progress in specific operational situations.
  \item \textbf{[R12] Paul Graham (2013), ``Do Things That Don't Scale.''} Classical practitioner essay on recruiting initial users manually to accelerate qualitative learning.
\end{itemize}

\subsection*{Additional Classical Texts Worth Studying}

\begin{itemize}[leftmargin=12pt,itemsep=2pt,topsep=1pt]
  \item \textbf{Carl Shapiro and Hal R. Varian, \textit{Information Rules} (1998):} Network externalities, customer lock-in, versioning, and switching costs.
  \item \textbf{Jean Tirole, \textit{The Theory of Industrial Organization} (1988):} Advanced formal mathematical models of pricing power, vertical integration, and oligopolistic bargaining.
  \item \textbf{Richard Rumelt, \textit{Good Strategy/Bad Strategy} (2011):} Dissecting strategic diagnosis, coherent guiding policies, and coordinated tactical action.
  \item \textbf{Andrew Chen, \textit{The Cold Start Problem} (2021):} Practical mechanics of overcoming network liquidity thresholds in two-sided platforms.
\end{itemize}

The foundational reading progression that best reinforces this chapter is: \textbf{Porter for industry structure $\to$ Gadiesh and Gilbert for profit-pool mapping $\to$ Christensen for jobs-to-be-done and disruptive asymmetries $\to$ Lieberman and Montgomery for entry timing $\to$ direct field discovery for your specific customer counterparty.}
